
\usepackage[T2A]{fontenc}
\usepackage[utf8]{inputenc}
\usepackage[russian,english]{babel}
\usepackage{amsmath,amssymb,amsfonts}
\usepackage[a4paper,margin=2.2cm]{geometry}
\usepackage{enumitem}
\setlist[enumerate,1]{label=\arabic*., itemsep=0.35em}
\setlist[enumerate,2]{label=\alph*), itemsep=0.15em}
\newcommand{\OPENNEGINF}{(-\infty}
\newcommand{\OPENPOSINF}{(+\infty}
\newcommand{\NEGINFTAIL}{-\infty}
\newcommand{\PLUSINFTAIL}{+\infty}
\title{Exam tasks (dec, jan, mar)}
\date{}
\begin{document}
\maketitle

\section*{Часть 1. Условия}

\subsection*{Множества — Операции над множествами}
\begin{enumerate}
\item \textbf{Декабрь, № 1} \par \smallskip
$\displaystyle \text{If the set A = }\{\text{1,2}\}\text{ and B = }\{\text{1,2,3}\}\text{, then which of the following is true?}$
\begin{enumerate}
\item $\displaystyle \text{A. B }\subseteq\text{ A}$
\item $\displaystyle \text{B. A=B}$
\item $\displaystyle \text{C. A}\subseteq\text{ B}$
\item $\displaystyle \text{D. A}\nsubseteq\text{ B}$
\end{enumerate}
\item \textbf{Январь, № 1} \par \smallskip
$\displaystyle \text{Let A = }\{\text{6,7,8,9,10}\}\text{ be a set, and }\emptyset\text{  denote the empty set. Then which of the following statements is correct?}$
\begin{enumerate}
\item $\displaystyle \text{A. 9 }\in\text{  A}$
\item $\displaystyle \text{B. \{6\} }\in\text{  A}$
\item $\displaystyle \text{C. 6 }\subseteq\text{  A}$
\item $\displaystyle \text{D. }\emptyset\text{ }\in\text{  A}$
\end{enumerate}
\item \textbf{Март, № 1} \par \smallskip
$\displaystyle \text{Let A = }\{\text{x }\mid\text{ -2 < x }\leq\text{  3}\}\text{ be a set. Which of the following statements is correct?}$
\begin{enumerate}
\item $\displaystyle \text{A. -3 }\in\text{  A}$
\item $\displaystyle \text{B. 3 }\in\text{  A}$
\item $\displaystyle \text{C. -2 }\in\text{  A}$
\item $\displaystyle \text{D. 2 }\notin\text{  A}$
\end{enumerate}
\item \textbf{Декабрь, № 2} \par \smallskip
$\displaystyle \text{If the set A = (-1, 3) and B = [2, 4], then what is A }\cap\text{  B?}$
\begin{enumerate}
\item $\displaystyle \text{A. [2, 3]}$
\item $\displaystyle \text{B. [2, 3)}$
\item $\displaystyle \text{C. (2, 3]}$
\item $\displaystyle \text{D. (2, 3)}$
\end{enumerate}
\item \textbf{Январь, № 2} \par \smallskip
$\displaystyle \text{Let A = }\{\text{x }\mid\text{ -2 }\leq\text{  x }\leq\text{  2}\}\text{ and B = }\{\text{x }\mid\text{ x > 0}\}\text{ be two sets. Then A }\cup\text{  B =}$
\begin{enumerate}
\item $\displaystyle \text{A. }\{\text{x }\mid\text{ x }\leq\text{  -2}\}$
\item $\displaystyle \text{B. }\{\text{x }\mid\text{ 0 < x }\leq\text{  2}\}$
\item $\displaystyle \text{C. }\{\text{x }\mid\text{ x }\geq\text{  -2}\}$
\item $\displaystyle \text{D. }\{\text{x }\mid\text{ -2 }\leq\text{  x < 0}\}$
\end{enumerate}
\item \textbf{Март, № 2} \par \smallskip
$\displaystyle \text{Let M = }\{\text{x }\mid\text{ -5 }\leq\text{  x < 2}\}\text{ , N = }\{\text{x }\mid\text{ -4 }\leq\text{  x }\leq\text{  2}\}\text{ be two sets. Then M }\cap\text{  N =}$
\begin{enumerate}
\item $\displaystyle \text{A. [-4,2)}$
\item $\displaystyle \text{B. (-5,2)}$
\item $\displaystyle \text{C. [-5,2]}$
\item $\displaystyle \text{D. (-4,2)}$
\end{enumerate}
\end{enumerate}

\subsection*{Неравенства — Квадратичные неравенства}
\begin{enumerate}
\item \textbf{Март, № 3} \par \smallskip
$\displaystyle \text{The solution set of the inequality x}^{2}\text{ - x - 6 > 0 is}$
\begin{enumerate}
\item $\displaystyle \text{A. }\{\text{x }\mid\text{ x < -2 or x > 3}\}$
\item $\displaystyle \text{B. }\{\text{x }\mid\text{ -2 < x < 3}\}$
\item $\displaystyle \text{C. }\{\text{x }\mid\text{ x > 3}\}$
\item $\displaystyle \text{D. }\{\text{x }\mid\text{ x < -2}\}$
\end{enumerate}
\item \textbf{Декабрь, № 3} \par \smallskip
$\displaystyle \text{The solution set of the inequality x}^{2}\text{ - 5x + 4 < 0 is}$
\begin{enumerate}
\item $\displaystyle \text{A. (-4, -1)}$
\item $\displaystyle \text{B. }\OPENNEGINF\text{ , 1) }\cup\text{ (4, }\PLUSINFTAIL\text{ )}$
\item $\displaystyle \text{C. }\OPENNEGINF\text{ , -4) }\cup\text{ (-1, }\PLUSINFTAIL\text{ )}$
\item $\displaystyle \text{D. (1, 4)}$
\end{enumerate}
\item \textbf{Январь, № 3} \par \smallskip
$\displaystyle \text{The solution set of the inequality x}^{2}\text{ - x - 2 > 0 is ( ).}$
\begin{enumerate}
\item $\displaystyle \text{A. }\{\text{x }\mid\text{ x < -1 or x > 2}\}$
\item $\displaystyle \text{B. }\{\text{x }\mid\text{ -1 < x < 2}\}$
\item $\displaystyle \text{C. }\{\text{x }\mid\text{ x < -2 or x > -1}\}$
\item $\displaystyle \text{D. }\{\text{x }\mid\text{ -2 < x }\leq\text{  1}\}$
\end{enumerate}
\end{enumerate}

\subsection*{Функции — Область определения}
\begin{enumerate}
\item \textbf{Март, № 4} \par \smallskip
$\displaystyle \text{The domain of the function f(x) = (1 - x}^{2}\text{)/(1 - x) is}$
\begin{enumerate}
\item $\displaystyle \text{A. }\OPENNEGINF\text{ ,0) }\cup\text{  (0, }\PLUSINFTAIL\text{ )}$
\item $\displaystyle \text{B. }\mathbb{R}$
\item $\displaystyle \text{C. }\OPENNEGINF\text{ , -1) }\cup\text{  (-1, }\PLUSINFTAIL\text{ )}$
\item $\displaystyle \text{D. }\OPENNEGINF\text{ ,1) }\cup\text{  (1, }\PLUSINFTAIL\text{ )}$
\end{enumerate}
\item \textbf{Декабрь, № 4} \par \smallskip
$\displaystyle \text{The domain of the function y = 1/}\sqrt{x+1}\text{ is}$
\begin{enumerate}
\item $\displaystyle \text{A. [-1, }\PLUSINFTAIL\text{ )}$
\item $\displaystyle \text{B. (-1, }\PLUSINFTAIL\text{ )}$
\item $\displaystyle \text{C. (0, }\PLUSINFTAIL\text{ )}$
\item $\displaystyle \text{D. }\OPENNEGINF\text{ , -1)}$
\end{enumerate}
\item \textbf{Январь, № 5} \par \smallskip
$\displaystyle \text{The domain of the function f(x) = 1/x + }\sqrt{1 - x}\text{ is ( ).}$
\begin{enumerate}
\item $\displaystyle \text{A. }\OPENNEGINF\text{ , 1]}$
\item $\displaystyle \text{B. }\OPENNEGINF\text{ , 0) }\cup\text{  (0, 1)}$
\item $\displaystyle \text{C. }\OPENNEGINF\text{ , 0) }\cup\text{  (0, 1]}$
\item $\displaystyle \text{D. [1, }\PLUSINFTAIL\text{ )}$
\end{enumerate}
\item \textbf{Декабрь, № 20} \par \smallskip
$\displaystyle \text{The domain of the function y = }\sqrt{x - 1}\text{ is}$
\begin{enumerate}
\item $\displaystyle \text{A. [1, }\PLUSINFTAIL\text{ )}$
\item $\displaystyle \text{B. (1, }\PLUSINFTAIL\text{ )}$
\item $\displaystyle \text{C. }\OPENNEGINF\text{ , 1]}$
\item $\displaystyle \text{D. }\OPENNEGINF\text{ , 1)}$
\end{enumerate}
\end{enumerate}

\subsection*{Функции — Обратные функции}
\begin{enumerate}
\item \textbf{Январь, № 9} \par \smallskip
$\displaystyle \text{The inverse function of the function y = x}^{3}\text{ + 3, x }\in\text{  }\mathbb{R}\text{ is ( ).}$
\begin{enumerate}
\item $\displaystyle \text{A. y = ∛(x - 3), x }\in\text{  }\mathbb{R}$
\item $\displaystyle \text{B. y = ∛(x + 3), x }\geq\text{  -3}$
\item $\displaystyle \text{C. y = ∛(x + 3), x }\in\text{  }\mathbb{R}$
\item $\displaystyle \text{D. y = ∛(x - 3), x }\geq\text{  3}$
\end{enumerate}
\item \textbf{Декабрь, № 9} \par \smallskip
$\displaystyle \text{The inverse function of y = x}^{3}\text{ is}$
\begin{enumerate}
\item $\displaystyle \text{A. y = x/3}$
\item $\displaystyle \text{B. y = 3 x}$
\item $\displaystyle \text{C. y = x}^{3}$
\item $\displaystyle \text{D. y = ∛x}$
\end{enumerate}
\item \textbf{Март, № 8} \par \smallskip
$\displaystyle \text{The inverse function of y = (3x - 2)/(2x + 1) is}$
\begin{enumerate}
\item $\displaystyle \text{A. y = (x + 2)/(2x - 3)}$
\item $\displaystyle \text{B. y = (2x - 3)/(x + 2)}$
\item $\displaystyle \text{C. y = -(2x - 3)/(x + 2)}$
\item $\displaystyle \text{D. y = -(x + 2)/(2x - 3)}$
\end{enumerate}
\end{enumerate}

\subsection*{Геометрия — Координатная геометрия}
\begin{enumerate}
\item \textbf{Январь, № 10} \par \smallskip
$\displaystyle \text{Which of the following points is in the third quadrant? ( ).}$
\begin{enumerate}
\item $\displaystyle \text{A. (-1,2)}$
\item $\displaystyle \text{B. (-1,-2)}$
\item $\displaystyle \text{C. (1,2)}$
\item $\displaystyle \text{D. (1,-2)}$
\end{enumerate}
\item \textbf{Декабрь, № 11} \par \smallskip
$\displaystyle \text{The point located in the third quadrant is}$
\begin{enumerate}
\item $\displaystyle \text{A. (-1, -2)}$
\item $\displaystyle \text{B. (-1, 2)}$
\item $\displaystyle \text{C. (2, 2)}$
\item $\displaystyle \text{D. (1, -2)}$
\end{enumerate}
\item \textbf{Март, № 5} \par \smallskip
$\displaystyle \text{If a = -2 , ab < 0 , then the point (a, b) belongs to}$
\begin{enumerate}
\item $\displaystyle \text{A. the third quadrant}$
\item $\displaystyle \text{B. the fourth quadrant}$
\item $\displaystyle \text{C. the second quadrant}$
\item $\displaystyle \text{D. the first quadrant}$
\end{enumerate}
\item \textbf{Декабрь, № 6} \par \smallskip
$\displaystyle \text{The point (-1, 3 - }\pi\text{ ) is located in the}$
\begin{enumerate}
\item $\displaystyle \text{A.  Third quadrant}$
\item $\displaystyle \text{B.  First quadrant}$
\item $\displaystyle \text{C.  Second quadrant}$
\item $\displaystyle \text{D.  Fourth quadrant}$
\end{enumerate}
\item \textbf{Январь, № 6} \par \smallskip
$\displaystyle \text{Let P(1,2) be a point in the rectangular coordinate system. If a point Q together with P are symmetrical about the x axis, then the coordinates of Q are ( ).}$
\begin{enumerate}
\item $\displaystyle \text{A. (2,1)}$
\item $\displaystyle \text{B. (1,-2)}$
\item $\displaystyle \text{C. (-1,2)}$
\item $\displaystyle \text{D. (-1,-2)}$
\end{enumerate}
\item \textbf{Март, № 16} \par \smallskip
$\displaystyle \text{Which of the following points is in the third quadrant?}$
\begin{enumerate}
\item $\displaystyle \text{A. (1,2)}$
\item $\displaystyle \text{B. (1,-2)}$
\item $\displaystyle \text{C. (-1,-2)}$
\item $\displaystyle \text{D. (-1,2)}$
\end{enumerate}
\end{enumerate}

\subsection*{Неравенства — Рациональные неравенства}
\begin{enumerate}
\item \textbf{Декабрь, № 12} \par \smallskip
$\displaystyle \text{The solution set of the inequality (x-1)/(x - 2) }\leq\text{  1  is}$
\begin{enumerate}
\item $\displaystyle \text{A. (2, }\PLUSINFTAIL\text{ )}$
\item $\displaystyle \text{B. [1, 2]}$
\item $\displaystyle \text{C. }\OPENNEGINF\text{ , 2)}$
\item $\displaystyle \text{D. }\OPENNEGINF\text{ , 1]}$
\end{enumerate}
\item \textbf{Март, № 11} \par \smallskip
$\displaystyle \text{The solution set of the inequality (2x - 3)/(3x - 4) }\leq\text{  0 is}$
\begin{enumerate}
\item $\displaystyle \text{A. [4/3, 3/2]}$
\item $\displaystyle \text{B. (4/3, 3/2)}$
\item $\displaystyle \text{C. (4/3, 3/2]}$
\item $\displaystyle \text{D. [4/3, 3/2)}$
\end{enumerate}
\item \textbf{Январь, № 12} \par \smallskip
$\displaystyle \text{The solution set of the rational inequality (2x + 1)/(x - 2) }\leq\text{  0 is ().}$
\begin{enumerate}
\item $\displaystyle \text{A. [-1/2, 2)}$
\item $\displaystyle \text{B. }\OPENNEGINF\text{ , -1/2] }\cup\text{  (2, }\PLUSINFTAIL\text{ )}$
\item $\displaystyle \text{C. [-1/2, 2]}$
\item $\displaystyle \text{D. }\OPENNEGINF\text{ , -1/2]}$
\end{enumerate}
\end{enumerate}

\subsection*{Среднее арифметическое и геометрическое — Среднее геометрическое}
\begin{enumerate}
\item \textbf{Март, № 12} \par \smallskip
$\displaystyle \text{The geometric mean of 3 and 12 is}$
\begin{enumerate}
\item $\displaystyle \text{A. ±10}$
\item $\displaystyle \text{B. ±8}$
\item $\displaystyle \text{C. ±4}$
\item $\displaystyle \text{D. ±6}$
\end{enumerate}
\end{enumerate}

\subsection*{Среднее арифметическое и геометрическое — Среднее арифметическое}
\begin{enumerate}
\item \textbf{Январь, № 13} \par \smallskip
$\displaystyle \text{Given that a = 2 - }\sqrt{3}\text{ and b = 2 + }\sqrt{3}\text{, the arithmetic mean of a and b is ().}$
\begin{enumerate}
\item $\displaystyle \text{A. 1}$
\item $\displaystyle \text{B. 2}$
\item $\displaystyle \text{C. ±1}$
\item $\displaystyle \text{D. -1}$
\end{enumerate}
\end{enumerate}

\subsection*{Функции — Равенство функций}
\begin{enumerate}
\item \textbf{Декабрь, № 15} \par \smallskip
$\displaystyle \text{Which of the following inequalities is correct?}$
\begin{enumerate}
\item $\displaystyle \text{A. 0.75^(-0.2) > 0.75^(-0.4)}$
\item $\displaystyle \text{B. 3^(4.5) > 3^5}$
\item $\displaystyle \text{C. 2.1^(-2) > 2.1^(-2)}$
\item $\displaystyle \text{D. 2.1^(2/3) > 1.2^(2/3)}$
\end{enumerate}
\end{enumerate}

\subsection*{Тригонометрия — Сторона угла через точку}
\begin{enumerate}
\item \textbf{Январь, № 15} \par \smallskip
$\displaystyle \text{If }\alpha\text{  is the angle between the x-axis and a line containing the point P(1,3), then }\sin\text{ }\alpha\text{  = ( ).}$
\begin{enumerate}
\item $\displaystyle \text{A. 3}$
\item $\displaystyle \text{B. 3}\sqrt{10}\text{/10}$
\item $\displaystyle \text{C. 1/3}$
\item $\displaystyle \text{D. }\sqrt{10}\text{/10}$
\end{enumerate}
\item \textbf{Декабрь, № 16} \par \smallskip
$\displaystyle \text{If point P(1,y) is on the terminal side of angle }\alpha\text{ , and }\tan\text{ }\alpha\text{  = 2, then y =}$
\begin{enumerate}
\item $\displaystyle \text{A. 1}$
\item $\displaystyle \text{B. 1/2}$
\item $\displaystyle \text{C. }\sqrt{3}$
\item $\displaystyle \text{D. 2}$
\end{enumerate}
\end{enumerate}

\subsection*{Тригонометрия — Табличные значения}
\begin{enumerate}
\item \textbf{Январь, № 7} \par \smallskip
$\displaystyle \text{Suppose that an angle }\alpha\text{  = 30^}\circ\text{ . Then which of the following statements is correct? ( ).}$
\begin{enumerate}
\item $\displaystyle \text{A. }\sin\text{ }\alpha\text{  = }\sqrt{2}\text{/2}$
\item $\displaystyle \text{B. }\cos\text{ }\alpha\text{  = }\sqrt{3}\text{/2}$
\item $\displaystyle \text{C. }\cos\text{ }\alpha\text{  = 1/2}$
\item $\displaystyle \text{D. }\tan\text{ }\alpha\text{  = }\sqrt{3}$
\end{enumerate}
\item \textbf{Декабрь, № 7} \par \smallskip
\textit{(no English statement.)}
\begin{enumerate}
\item $\displaystyle \text{A. 2}\sqrt{2}$
\item $\displaystyle \text{B. 2}$
\item $\displaystyle \text{C. }\sqrt{2}$
\item $\displaystyle \text{D. 1}$
\end{enumerate}
\item \textbf{Декабрь, № 22} \par \smallskip
$\displaystyle \text{If }\alpha\text{  is an acute angle, and }\sin\text{ }\alpha\text{  = }\sqrt{3}\text{/2, then }\tan\text{ }\alpha\text{  =}$
\begin{enumerate}
\item $\displaystyle \text{A. 1/2}$
\item $\displaystyle \text{B. }\sqrt{3}$
\item $\displaystyle \text{C. 1/4}$
\item $\displaystyle \text{D. }\sqrt{3}\text{/3}$
\end{enumerate}
\item \textbf{Март, № 6} \par \smallskip
$\displaystyle \text{Let }\alpha\text{  = (5/6)}\pi\text{ . Then }\cos\text{ }\alpha\text{  =}$
\begin{enumerate}
\item $\displaystyle \text{A. -}\sqrt{3}\text{/2}$
\item $\displaystyle \text{B. }\sqrt{3}\text{/2}$
\item $\displaystyle \text{C. -1/2}$
\item $\displaystyle \text{D. 1/2}$
\end{enumerate}
\item \textbf{Январь, № 20} \par \smallskip
$\displaystyle \text{Given that }\sin\text{ }\alpha\text{  = 4/5, and the angle }\alpha\text{  belongs to the second quadrant, then which of the following statements is correct? ( )}$
\begin{enumerate}
\item $\displaystyle \text{A. }\cos\text{ }\alpha\text{  = 3/5}$
\item $\displaystyle \text{B. }\cos\text{ }\alpha\text{  = -3/5}$
\item $\displaystyle \text{C. }\tan\text{ }\alpha\text{  = -3/4}$
\item $\displaystyle \text{D. }\tan\text{ }\alpha\text{  = 4/3}$
\end{enumerate}
\item \textbf{Декабрь, № 34} \par \smallskip
$\displaystyle \text{Given that }\cos\text{ a = -1/2, }\alpha\text{ }\in\text{ (}\pi\text{ /2, }\pi\text{ ), then the value of }\sin\text{ }\alpha\text{  is}$
\begin{enumerate}
\item $\displaystyle \text{A. -}\sqrt{3}\text{/2}$
\item $\displaystyle \text{B. }\sqrt{3}\text{/2}$
\item $\displaystyle \text{C. 1/2}$
\item $\displaystyle \text{D. -1/2}$
\end{enumerate}
\item \textbf{Декабрь, № 19} \par \smallskip
$\displaystyle \text{Given that }\sin\text{ }\alpha\text{  = 3/5, and }\alpha\text{  is in the second quadrant, then }\cos\text{ }\alpha\text{  =}$
\begin{enumerate}
\item $\displaystyle \text{A. 4/5}$
\item $\displaystyle \text{B. -4/5}$
\item $\displaystyle \text{C. 3/4}$
\item $\displaystyle \text{D. -3/4}$
\end{enumerate}
\end{enumerate}

\subsection*{Функции — Свойства функций}
\begin{enumerate}
\item \textbf{Март, № 7} \par \smallskip
$\displaystyle \text{Over its domain, the function f(x) = }\ln\text{ (x}^{2}\text{) is}$
\begin{enumerate}
\item $\displaystyle \text{A. uncertain about its parity}$
\item $\displaystyle \text{B. neither odd nor even}$
\item $\displaystyle \text{C. odd}$
\item $\displaystyle \text{D. even}$
\end{enumerate}
\item \textbf{Январь, № 8} \par \smallskip
$\displaystyle \text{About the function f(x) = x}^{4}\text{ + 3, which of the following statements is correct? ( ).}$
\begin{enumerate}
\item $\displaystyle \text{A.  / The function is odd}$
\item $\displaystyle \text{B.  / None is true}$
\item $\displaystyle \text{C.  / Neither odd nor even}$
\item $\displaystyle \text{D.  / The function is even}$
\end{enumerate}
\item \textbf{Март, № 9} \par \smallskip
$\displaystyle \text{The function f(x) = }\sin\text{ x is monotonic increasing over}$
\begin{enumerate}
\item $\displaystyle \text{A. [-}\pi\text{ /2, }\pi\text{ /2]}$
\item $\displaystyle \text{B. [-}\pi\text{ , }\pi\text{ ]}$
\item $\displaystyle \text{C. }\OPENNEGINF\text{ , }\PLUSINFTAIL\text{ )}$
\item $\displaystyle \text{D. [0, }\pi\text{ ]}$
\end{enumerate}
\item \textbf{Январь, № 11} \par \smallskip
$\displaystyle \text{Given that y = |x|, which of the following conclusions is correct? ().}$
\begin{enumerate}
\item $\displaystyle \text{A.  x < 0  When x<0 it's an increasing function}$
\item $\displaystyle \text{B.  x > 0  When x>0 it's an increasing function}$
\item $\displaystyle \text{C.  x }\in\text{  }\mathbb{R}\text{  When x}\in\text{ R is's a decreasing function}$
\item $\displaystyle \text{D.  x }\in\text{  }\mathbb{R}\text{  When x}\in\text{ R is's a increasing function}$
\end{enumerate}
\item \textbf{Март, № 31} \par \smallskip
$\displaystyle \text{The minimum positive period of the function f(x) = 3cos(4x - }\pi\text{ /3) is}$
\begin{enumerate}
\item $\displaystyle \text{A. }\pi\text{ /4}$
\item $\displaystyle \text{B. 2}\pi\text{ }$
\item $\displaystyle \text{C. }\pi\text{ /2}$
\item $\displaystyle \text{D. }\pi\text{ /3}$
\end{enumerate}
\item \textbf{Декабрь, № 8} \par \smallskip
$\displaystyle \text{Which of the following functions is an odd function?}$
\begin{enumerate}
\item $\displaystyle \text{A. y = x }\sin\text{ x}$
\item $\displaystyle \text{B. y = x^4}$
\item $\displaystyle \text{C. y = }\cos\text{ x}$
\item $\displaystyle \text{D. y = x + x}^{3}$
\end{enumerate}
\item \textbf{Декабрь, № 10} \par \smallskip
$\displaystyle \text{Which of the following functions is monotonically increasing?}$
\begin{enumerate}
\item $\displaystyle \text{A. y = (1/3)^x}$
\item $\displaystyle \text{B. y = log\_0.5 x}$
\item $\displaystyle \text{C. y = 2x + 1}$
\item $\displaystyle \text{D. y = }\sin\text{ x}$
\end{enumerate}
\item \textbf{Январь, № 23} \par \smallskip
$\displaystyle \text{About the exponential function y = aˣ (a > 0, and a }\neq\text{  1), which of the following statements is incorrect? ( )}$
\begin{enumerate}
\item $\displaystyle \text{A.   The range is (0, }\PLUSINFTAIL\text{ )}$
\item $\displaystyle \text{B.  Graph passes through (0, 1)}$
\item $\displaystyle \text{C.   It's increasing}$
\item $\displaystyle \text{D.  The domain is }\mathbb{R}$
\end{enumerate}
\item \textbf{Декабрь, № 26} \par \smallskip
$\displaystyle \text{Which of the following descriptions about the function y = (1/2)^x is correct?}$
\begin{enumerate}
\item $\displaystyle \text{A.  Symmetric about the origin}$
\item $\displaystyle \text{B. x Symmetric about the x-axis}$
\item $\displaystyle \text{C.  Increasing}$
\item $\displaystyle \text{D.  Decreasing}$
\end{enumerate}
\item \textbf{Март, № 36} \par \smallskip
$\displaystyle \text{The range of the function f(x) = 1/(1 + |x|) (x }\in\text{  }\mathbb{R}\text{) is}$
\begin{enumerate}
\item $\displaystyle \text{A. (0, 1]}$
\item $\displaystyle \text{B. [0, 1]}$
\item $\displaystyle \text{C. [0, 1)}$
\item $\displaystyle \text{D. (0, 1)}$
\end{enumerate}
\item \textbf{Март, № 25} \par \smallskip
$\displaystyle \text{If an exponential function y = aˣ is monotonic increasing on }\mathbb{R}\text{, then the range of a is}$
\begin{enumerate}
\item $\displaystyle \text{A. a > 1}$
\item $\displaystyle \text{B. a < 0 or a > 1}$
\item $\displaystyle \text{C. 0 < a < 1}$
\item $\displaystyle \text{D. a > 0}$
\end{enumerate}
\item \textbf{Март, № 40} \par \smallskip
$\displaystyle \text{Which of the following descriptions about the function f(x) = }\tan\text{ x is incorrect?}$
\begin{enumerate}
\item $\displaystyle \text{A. f(x) satisfies f(-x) = -f(x) on its domain}$
\item $\displaystyle \text{B. f(x) is even}$
\item $\displaystyle \text{C. f(x) is monotonic increasing over (-}\pi\text{ /2, }\pi\text{ /2)}$
\item $\displaystyle \text{D. f(x) has exactly one zero on the interval (-}\pi\text{ /2, }\pi\text{ /2)}$
\end{enumerate}
\end{enumerate}

\subsection*{Последовательности — Простые задачи}
\begin{enumerate}
\item \textbf{Март, № 10} \par \smallskip
$\displaystyle \text{Let \{a}_{n}\text{\} be an arithmetic sequence, with the first term a}_{1}\text{ = 2 and common difference d = 3. Then a}_{100}\text{ =}$
\begin{enumerate}
\item $\displaystyle \text{A. 298}$
\item $\displaystyle \text{B. 299}$
\item $\displaystyle \text{C. 302}$
\item $\displaystyle \text{D. 305}$
\end{enumerate}
\item \textbf{Март, № 22} \par \smallskip
$\displaystyle \text{In an arithmetic sequence \{a}_{n}\text{\} , if a}_{1}\text{ = 1 and its common difference d = 2 , then a}_{10}\text{ =}$
\begin{enumerate}
\item $\displaystyle \text{A. 19}$
\item $\displaystyle \text{B. 17}$
\item $\displaystyle \text{C. 15}$
\item $\displaystyle \text{D. 21}$
\end{enumerate}
\item \textbf{Январь, № 4} \par \smallskip
$\displaystyle \text{Let \{a}_{n}\text{\} be an arithmetic sequence, with the first term a}_{1}\text{ = 2 and common difference d = 3. Then a}_{100}\text{ = ( ).}$
\begin{enumerate}
\item $\displaystyle \text{A. 299}$
\item $\displaystyle \text{B. 302}$
\item $\displaystyle \text{C. 305}$
\item $\displaystyle \text{D. 298}$
\end{enumerate}
\item \textbf{Декабрь, № 5} \par \smallskip
$\displaystyle \text{If \{a}_{n}\text{\} is an arithmetic sequence, the first term a}_{1}\text{ = 5, and d = 2, then a}_{100}\text{ =}$
\begin{enumerate}
\item $\displaystyle \text{A. 203}$
\item $\displaystyle \text{B. 205}$
\item $\displaystyle \text{C. 199}$
\item $\displaystyle \text{D. 204}$
\end{enumerate}
\item \textbf{Декабрь, № 13} \par \smallskip
$\displaystyle \text{Given that a, b and c form an arithmetic sequence and satisfy a + c = 20, then b =}$
\begin{enumerate}
\item $\displaystyle \text{A. 10}$
\item $\displaystyle \text{B. 5}$
\item $\displaystyle \text{C. 20}$
\item $\displaystyle \text{D. 0}$
\end{enumerate}
\item \textbf{Январь, № 16} \par \smallskip
$\displaystyle \text{Among the four sequences below, the number of arithmetic sequences is ( ).  
  (1) 7,13,19,25  (2) 2,4,7,11   (3) -1,-3,-5,-7    (4) 2,4,8,16}$
\begin{enumerate}
\item $\displaystyle \text{A. 2}$
\item $\displaystyle \text{B. 3}$
\item $\displaystyle \text{C. 1}$
\item $\displaystyle \text{D. 4}$
\end{enumerate}
\item \textbf{Декабрь, № 17} \par \smallskip
$\displaystyle \text{Given that \{a}_{n}\text{\} is an arithmetic sequence, and a}_{1}\text{ = 1, a}_{2}\text{ = 2, then a}_{2025}\text{ =}$
\begin{enumerate}
\item $\displaystyle \text{A. 2}$
\item $\displaystyle \text{B. 2025}$
\item $\displaystyle \text{C. 2024}$
\item $\displaystyle \text{D. 20}$
\end{enumerate}
\item \textbf{Январь, № 21} \par \smallskip
$\displaystyle \text{In an arithmetic sequence \{a}_{n}\text{\}, if a}_{2}\text{ = 1 and a}_{4}\text{ = 5, then the first term a}_{1}\text{ and common difference d are ( ) respectively.}$
\begin{enumerate}
\item $\displaystyle \text{A. -1,2}$
\item $\displaystyle \text{B. 1, -2}$
\item $\displaystyle \text{C. -1, -2}$
\item $\displaystyle \text{D. 1,2}$
\end{enumerate}
\item \textbf{Март, № 15} \par \smallskip
$\displaystyle \text{In a geometric sequence \{a}_{n}\text{\} , if a}_{1}\text{ = 4 and a}_{4}\text{ = 32 , then a}_{3}\text{ is}$
\begin{enumerate}
\item $\displaystyle \text{A. 8}$
\item $\displaystyle \text{B. 16}$
\item $\displaystyle \text{C. -8}$
\item $\displaystyle \text{D. -16}$
\end{enumerate}
\item \textbf{Декабрь, № 23} \par \smallskip
$\displaystyle \text{Given that the first three terms of the geometric sequence \{a}_{n}\text{\} are 1, 2, 4, then the general term formula of the sequence is}$
\begin{enumerate}
\item $\displaystyle \text{A. a}_{n}\text{ = 2n}$
\item $\displaystyle \text{B. a}_{n}\text{ = 8}$
\item $\displaystyle \text{C. a}_{n}\text{ = 2^(n-1)}$
\item $\displaystyle \text{D. a}_{n}\text{ = 2^n}$
\end{enumerate}
\item \textbf{Январь, № 40} \par \smallskip
$\displaystyle \text{The general formula for the sequence -1/5, 1/7, -1/9, 1/11 is a}_{n}\text{ = ( ).}$
\begin{enumerate}
\item $\displaystyle \text{A. (-1)^(n-1) / (3n + 2)}$
\item $\displaystyle \text{B. (-1)^n / (2n + 3)}$
\item $\displaystyle \text{C. (-1)^(n-1) / (2n + 3)}$
\item $\displaystyle \text{D. (-1)^n / (3n + 2)}$
\end{enumerate}
\item \textbf{Март, № 41} \par \smallskip
$\displaystyle \text{In a sequence \{a}_{n}\text{\} , if a}_{1}\text{ = 1/2 and 1/a}_{n}\text{₊₁ = 1/a}_{n}\text{ + 1, n }\geq\text{  1, then a}_{6}\text{ =}$
\begin{enumerate}
\item $\displaystyle \text{A. 1/7}$
\item $\displaystyle \text{B. 1/6}$
\item $\displaystyle \text{C. 1/9}$
\item $\displaystyle \text{D. 1/8}$
\end{enumerate}
\item \textbf{Март, № 47} \par \smallskip
$\displaystyle \text{Suppose that a}_{n}\text{ = 2}^{n}\text{⁺¹ - 3n, n }\geq\text{  1. Let S}_{n}\text{ be the sum of the first n terms of the sequence. Then S}_{n}\text{ =}$
\begin{enumerate}
\item $\displaystyle \text{A. 2}^{n}\text{⁺² - 3n(n+1)/2 + 4}$
\item $\displaystyle \text{B. 2}^{n}\text{⁺² - 3n - 4}$
\item $\displaystyle \text{C. 2}^{n}\text{⁺² + 3n(n+1)/2 - 4}$
\item $\displaystyle \text{D. 2}^{n}\text{⁺² - 3n(n+1)/2 - 4}$
\end{enumerate}
\end{enumerate}

\subsection*{Неравенства — Свойства неравенств}
\begin{enumerate}
\item \textbf{Январь, № 22} \par \smallskip
$\displaystyle \text{Let a, b, c be three real numbers. Which of the following statements is correct? ( )}$
\begin{enumerate}
\item $\displaystyle \text{A.  If a > b then ac > bc}$
\item $\displaystyle \text{B.  If a > b then a + c > b + c}$
\item $\displaystyle \text{C.  If ac > bc then a > b}$
\item $\displaystyle \text{D.  If a > b then a}^{2}\text{ > b}^{2}$
\end{enumerate}
\item \textbf{Март, № 23} \par \smallskip
$\displaystyle \text{Let a > b. Then which of the following is correct?}$
\begin{enumerate}
\item $\displaystyle \text{A. a}^{3}\text{ > b}^{3}$
\item $\displaystyle \text{B. |a| > |b|}$
\item $\displaystyle \text{C. a}^{2}\text{ > b}^{2}$
\item $\displaystyle \text{D. |a| < |b|}$
\end{enumerate}
\item \textbf{Март, № 19} \par \smallskip
$\displaystyle \text{Which of the following inequalities is correct?}$
\begin{enumerate}
\item $\displaystyle \text{A. 2.1^(2/3) > 1.2^(2/3)}$
\item $\displaystyle \text{B. 2.1^(-2) > 1.2^(-2)}$
\item $\displaystyle \text{C. 3^(2.25) > 3^3}$
\item $\displaystyle \text{D. 0.75^(-0.2) > 0.75^(-0.4)}$
\end{enumerate}
\end{enumerate}

\subsection*{Конические сечения — Окружность}
\begin{enumerate}
\item \textbf{Март, № 20} \par \smallskip
$\displaystyle \text{A circle centered at A(-3, 2), with a radius 3 has equation}$
\begin{enumerate}
\item $\displaystyle \text{A. (x - 3)}^{2}\text{ + (y + 2)}^{2}\text{ = 3}$
\item $\displaystyle \text{B. (x + 3)}^{2}\text{ + (y - 2)}^{2}\text{ = 9}$
\item $\displaystyle \text{C. (x - 3)}^{2}\text{ + (y + 2)}^{2}\text{ = 9}$
\item $\displaystyle \text{D. (x + 3)}^{2}\text{ + (y - 2)}^{2}\text{ = 3}$
\end{enumerate}
\item \textbf{Март, № 24} \par \smallskip
$\displaystyle \text{Which of the following is the equation of a circle with center (-3,2) and radius 2?}$
\begin{enumerate}
\item $\displaystyle \text{A. (x + 3)}^{2}\text{ + (y - 2)}^{2}\text{ = 16}$
\item $\displaystyle \text{B. (x + 3)}^{2}\text{ + (y - 2)}^{2}\text{ = 4}$
\item $\displaystyle \text{C. (x - 3)}^{2}\text{ + (y + 2)}^{2}\text{ = 4}$
\item $\displaystyle \text{D. (x - 3)}^{2}\text{ + (y + 2)}^{2}\text{ = 16}$
\end{enumerate}
\item \textbf{Январь, № 28} \par \smallskip
$\displaystyle \text{A circle has center at (-3,2) and radius 4. Then the equation of the circle is ( ).}$
\begin{enumerate}
\item $\displaystyle \text{A. (x + 3)}^{2}\text{ + (y - 2)}^{2}\text{ = 4}$
\item $\displaystyle \text{B. (x - 3)}^{2}\text{ + (y + 2)}^{2}\text{ = 4}$
\item $\displaystyle \text{C. (x - 3)}^{2}\text{ + (y + 2)}^{2}\text{ = 16}$
\item $\displaystyle \text{D. (x + 3)}^{2}\text{ + (y - 2)}^{2}\text{ = 16}$
\end{enumerate}
\item \textbf{Декабрь, № 18} \par \smallskip
$\displaystyle \text{If the center of the circle is (-3,2) and the radius is 4, then the equation of the circle is}$
\begin{enumerate}
\item $\displaystyle \text{A. (x + 3)}^{2}\text{ + (y - 2)}^{2}\text{= 4}$
\item $\displaystyle \text{B. (x - 3)}^{2}\text{ + (y + 2)}^{2}\text{ = 4}$
\item $\displaystyle \text{C. (x - 3)}^{2}\text{ + (y + 2)}^{2}\text{ = 16}$
\item $\displaystyle \text{D. (x + 3)}^{2}\text{ + (y - 2)}^{2}\text{ = 16}$
\end{enumerate}
\item \textbf{Январь, № 19} \par \smallskip
$\displaystyle \text{If the equation of a circle is x}^{2}\text{ + y}^{2}\text{ - 4x - 3 = 0, then the center and radius are given by ( ).}$
\begin{enumerate}
\item $\displaystyle \text{A. (0,2)  }\sqrt{7}$
\item $\displaystyle \text{B. (0,2)  7}$
\item $\displaystyle \text{C. (2,0)  }\sqrt{7}$
\item $\displaystyle \text{D. (2,0)  7}$
\end{enumerate}
\item \textbf{Декабрь, № 21} \par \smallskip
$\displaystyle \text{Given the equation of the circle is x}^{2}\text{ + y}^{2}\text{ + 2x - 2y = 0, then the radius of the circle is}$
\begin{enumerate}
\item $\displaystyle \text{A. 2}$
\item $\displaystyle \text{B. 1/2}$
\item $\displaystyle \text{C. }\sqrt{2}$
\item $\displaystyle \text{D. 1}$
\end{enumerate}
\end{enumerate}

\subsection*{Геометрия — Формула расстояния}
\begin{enumerate}
\item \textbf{Январь, № 17} \par \smallskip
$\displaystyle \text{The distance from a point A(3, -2) to B(-5, -1) on the plane is |AB| = ( ).}$
\begin{enumerate}
\item $\displaystyle \text{A. 65}$
\item $\displaystyle \text{B. }\sqrt{65}$
\item $\displaystyle \text{C. 13}$
\item $\displaystyle \text{D. }\sqrt{13}$
\end{enumerate}
\item \textbf{Декабрь, № 25} \par \smallskip
$\displaystyle \text{The distance from point P(-1, 2) to point Q(3, 1) is}$
\begin{enumerate}
\item $\displaystyle \text{A. }\sqrt{17}$
\item $\displaystyle \text{B. }\sqrt{13}$
\item $\displaystyle \text{C. }\sqrt{5}$
\item $\displaystyle \text{D. 5}$
\end{enumerate}
\item \textbf{Март, № 30} \par \smallskip
$\displaystyle \text{The distance between the points P(-1,2) and Q(3,1) is}$
\begin{enumerate}
\item $\displaystyle \text{A. }\sqrt{5}$
\item $\displaystyle \text{B. 5}$
\item $\displaystyle \text{C. }\sqrt{13}$
\item $\displaystyle \text{D. }\sqrt{17}$
\end{enumerate}
\item \textbf{Январь, № 33} \par \smallskip
$\displaystyle \text{The distance between the points P(-1,2) and Q(3,1) is ( ).}$
\begin{enumerate}
\item $\displaystyle \text{A. }\sqrt{5}$
\item $\displaystyle \text{B. }\sqrt{17}$
\item $\displaystyle \text{C. }\sqrt{13}$
\item $\displaystyle \text{D. 5}$
\end{enumerate}
\item \textbf{Март, № 17} \par \smallskip
$\displaystyle \text{Let A(-2, -1) and B(a,3) be two points, and |AB| = 5. Then the value of a is}$
\begin{enumerate}
\item $\displaystyle \text{A. ±5}$
\item $\displaystyle \text{B. ±1}$
\item $\displaystyle \text{C. -1 or 5}$
\item $\displaystyle \text{D. 1 or -5}$
\end{enumerate}
\end{enumerate}

\subsection*{Функции — Неравенства}
\begin{enumerate}
\item \textbf{Декабрь, № 24} \par \smallskip
$\displaystyle \text{Given that a > b, then}$
\begin{enumerate}
\item $\displaystyle \text{A. 1/a < 1/b}$
\item $\displaystyle \text{B. a^2 > b^2}$
\item $\displaystyle \text{C. 5^a > 5^b}$
\item $\displaystyle \text{D. }\sin\text{ a > }\sin\text{ b}$
\end{enumerate}
\item \textbf{Январь, № 27} \par \smallskip
$\displaystyle \text{Which of the following inequalities is correct? ( )}$
\begin{enumerate}
\item $\displaystyle \text{A. 2.1⁻² > 1.2⁻²}$
\item $\displaystyle \text{B. 0.75⁻⁰}\cdot\text{ ² > 0.75⁻⁰}\cdot\text{ ⁴}$
\item $\displaystyle \text{C. 3}^{2}\cdot\text{ ²}^{5}\text{ > 3}^{3}$
\item $\displaystyle \text{D. 2.1^(2/3) > 1.2^(2/3)}$
\end{enumerate}
\end{enumerate}

\subsection*{Тригонометрия — Формулы двойного угла}
\begin{enumerate}
\item \textbf{Март, № 26} \par \smallskip
$\displaystyle \text{If }\sin\text{ a = }\sqrt{3}\text{/3, then }\cos\text{ 2a =}$
\begin{enumerate}
\item $\displaystyle \text{A. 2/3}$
\item $\displaystyle \text{B. 1/3}$
\item $\displaystyle \text{C. -1/3}$
\item $\displaystyle \text{D. -2/3}$
\end{enumerate}
\item \textbf{Декабрь, № 27} \par \smallskip
$\displaystyle \text{If }\sin\text{ }\alpha\text{  = 1/4, then }\cos\text{ 2}\alpha\text{  =}$
\begin{enumerate}
\item $\displaystyle \text{A. -7/8}$
\item $\displaystyle \text{B. 1/8}$
\item $\displaystyle \text{C. 7/8}$
\item $\displaystyle \text{D. -1/8}$
\end{enumerate}
\item \textbf{Январь, № 24} \par \smallskip
$\displaystyle \text{If }\sin\text{ }\alpha\text{  = 3/5, and }\alpha\text{  belongs to the first quadrant, then }\sin\text{ 2}\alpha\text{  = ( ).}$
\begin{enumerate}
\item $\displaystyle \text{A. 12/25}$
\item $\displaystyle \text{B. 24/25}$
\item $\displaystyle \text{C. 18/25}$
\item $\displaystyle \text{D. 7/25}$
\end{enumerate}
\end{enumerate}

\subsection*{Тригонометрия — Тригонометрические выражения}
\begin{enumerate}
\item \textbf{Март, № 14} \par \smallskip
$\displaystyle \text{If a triangle ABC has sides BC = 5, AC = 12, AB = 13, then }\sin\text{ A is}$
\begin{enumerate}
\item $\displaystyle \text{A. 5/13}$
\item $\displaystyle \text{B. 5/12}$
\item $\displaystyle \text{C. 12/13}$
\item $\displaystyle \text{D. 7/12}$
\end{enumerate}
\item \textbf{Декабрь, № 40} \par \smallskip
$\displaystyle \text{Given that (}\cos\text{ a - }\sin\text{ a)/(}\cos\text{ a + }\sin\text{ a) = 1/3, then }\tan\text{ a =}$
\begin{enumerate}
\item $\displaystyle \text{A. 1/2}$
\item $\displaystyle \text{B. 1/4}$
\item $\displaystyle \text{C. 2/3}$
\item $\displaystyle \text{D. 1/3}$
\end{enumerate}
\end{enumerate}

\subsection*{Конические сечения — Гипербола}
\begin{enumerate}
\item \textbf{Март, № 28} \par \smallskip
$\displaystyle \text{Let mx}^{2}\text{ + ny}^{2}\text{ = 1 be a hyperbola. Then m, n satisfy}$
\begin{enumerate}
\item $\displaystyle \text{A. mn > 1}$
\item $\displaystyle \text{B. mn < 0}$
\item $\displaystyle \text{C. mn < 1}$
\item $\displaystyle \text{D. mn > 0}$
\end{enumerate}
\item \textbf{Декабрь, № 29} \par \smallskip
$\displaystyle \text{The coordinates of the foci of the hyperbola x}^{2}\text{ - y}^{2}\text{ = 1 are}$
\begin{enumerate}
\item $\displaystyle \text{A. (±}\sqrt{2}\text{, 0)}$
\item $\displaystyle \text{B. (0, ±2)}$
\item $\displaystyle \text{C. (0, ±}\sqrt{2}\text{)}$
\item $\displaystyle \text{D. (±2, 0)}$
\end{enumerate}
\item \textbf{Январь, № 26} \par \smallskip
$\displaystyle \text{Given that the equation of a hyperbola is x}^{2}\text{/64 - y}^{2}\text{/16 = 1, which of the following statements is correct? ( )}$
\begin{enumerate}
\item $\displaystyle \text{A.  / focal distance = 8}\sqrt{5}$
\item $\displaystyle \text{B.  / semi-major axis  = 4}$
\item $\displaystyle \text{C.  /eccentricity  = }\sqrt{5}$
\item $\displaystyle \text{D.  / semi-minor axis = 8}$
\end{enumerate}
\end{enumerate}

\subsection*{Геометрия — Прямые}
\begin{enumerate}
\item \textbf{Март, № 13} \par \smallskip
$\displaystyle \text{The angle of inclination of the line x + }\sqrt{3}\text{ y - 2 = 0 is}$
\begin{enumerate}
\item $\displaystyle \text{A. 150^}\circ\text{ }$
\item $\displaystyle \text{B. 60^}\circ\text{ }$
\item $\displaystyle \text{C. 30^}\circ\text{ }$
\item $\displaystyle \text{D. 120^}\circ\text{ }$
\end{enumerate}
\item \textbf{Декабрь, № 14} \par \smallskip
$\displaystyle \text{The inclination angle of the line y = }\sqrt{3}\text{x + 10 is}$
\begin{enumerate}
\item $\displaystyle \text{A. 60^}\circ\text{ }$
\item $\displaystyle \text{B. 30^}\circ\text{ }$
\item $\displaystyle \text{C. 45^}\circ\text{ }$
\item $\displaystyle \text{D. 120^}\circ\text{ }$
\end{enumerate}
\item \textbf{Март, № 18} \par \smallskip
$\displaystyle \text{Suppose that a line l passes through two points M(0, -1) and N(-2, 0). Then the equation of l is}$
\begin{enumerate}
\item $\displaystyle \text{A. 2x + y + 2 = 0}$
\item $\displaystyle \text{B. x + y + 2 = 0}$
\item $\displaystyle \text{C. 2x + y + 1 = 0}$
\item $\displaystyle \text{D. x + 2y + 2 = 0}$
\end{enumerate}
\item \textbf{Март, № 33} \par \smallskip
$\displaystyle \text{The line that is parallel to 4x - 2y - 1 = 0 is}$
\begin{enumerate}
\item $\displaystyle \text{A. x + 2y - 1 = 0}$
\item $\displaystyle \text{B. 2x - y - 1 = 0}$
\item $\displaystyle \text{C. 2x + y - 1 = 0}$
\item $\displaystyle \text{D. x - 2y - 1 = 0}$
\end{enumerate}
\item \textbf{Январь, № 14} \par \smallskip
$\displaystyle \text{If the straight line l passes through the points A(-2,3) and B(3,1), then the slope of l is ().}$
\begin{enumerate}
\item $\displaystyle \text{A. 2/5}$
\item $\displaystyle \text{B. 5/2}$
\item $\displaystyle \text{C. -2/5}$
\item $\displaystyle \text{D. -5/2}$
\end{enumerate}
\item \textbf{Январь, № 18} \par \smallskip
$\displaystyle \text{If a straight line l has an angle of inclination of 45^}\circ\text{ , and it passes through the point (0,2), then the equation of l is ( ).}$
\begin{enumerate}
\item $\displaystyle \text{A. y = 2x - 2}$
\item $\displaystyle \text{B. y = x + 2}$
\item $\displaystyle \text{C. y = 2x + 2}$
\item $\displaystyle \text{D. y = x - 2}$
\end{enumerate}
\item \textbf{Январь, № 25} \par \smallskip
$\displaystyle \text{The intersection of the lines l}_{1}\text{: 3x - y + 8 = 0 and l}_{2}\text{: x + 2y - 9 = 0 has coordinates ( ).}$
\begin{enumerate}
\item $\displaystyle \text{A. (-5,1)}$
\item $\displaystyle \text{B. (1,-5)}$
\item $\displaystyle \text{C. (5,-1)}$
\item $\displaystyle \text{D. (-1,5)}$
\end{enumerate}
\item \textbf{Март, № 27} \par \smallskip
$\displaystyle \text{If three lines ax + 2y + 8 = 0, 4x + 3y = 10 and 2x - y = 10 intersect at one point, then a =}$
\begin{enumerate}
\item $\displaystyle \text{A. 1}$
\item $\displaystyle \text{B. -1}$
\item $\displaystyle \text{C. 3}$
\item $\displaystyle \text{D. -3}$
\end{enumerate}
\item \textbf{Декабрь, № 28} \par \smallskip
$\displaystyle \text{The intersection point of line l}_{1}\text{: y = 2x + 1 and line l}_{2}\text{: x + y + 1 = 0 is}$
\begin{enumerate}
\item $\displaystyle \text{A. (1/3, 2/3)}$
\item $\displaystyle \text{B. (2/3, 1/3)}$
\item $\displaystyle \text{C. (-1/3, -2/3)}$
\item $\displaystyle \text{D. (-2/3, -1/3)}$
\end{enumerate}
\item \textbf{Декабрь, № 33} \par \smallskip
$\displaystyle \text{Which of the following lines is perpendicular to the line l: x + y + 3 = 0?}$
\begin{enumerate}
\item $\displaystyle \text{A. y = -2x + 1}$
\item $\displaystyle \text{B. y = x + 2}$
\item $\displaystyle \text{C. y = 1 - x}$
\item $\displaystyle \text{D. y = 2x + 2}$
\end{enumerate}
\item \textbf{Декабрь, № 44} \par \smallskip
$\displaystyle \text{If line l is perpendicular to line l}_{1}\text{:x + 2y + 4 = 0 and passes through the intersection point of lines l}_{2}\text{:x + y + 1 = 0 and l}_{3}\text{:2x + y - 1 = 0, then the equation of l is}$
\begin{enumerate}
\item $\displaystyle \text{A. x + 2y - 3 = 0}$
\item $\displaystyle \text{B. 2x + y - 5 = 0}$
\item $\displaystyle \text{C. 2x - y - 7 = 0}$
\item $\displaystyle \text{D. x - 2y - 1 = 0}$
\end{enumerate}
\item \textbf{Январь, № 45} \par \smallskip
$\displaystyle \text{Suppose that a straight line l is perpendicular to another line l}_{1}\text{: x + 2y + 4 = 0 and passes through the intersection of l}_{2}\text{: x + y + 1 = 0 and l}_{3}\text{: 2x + y - 1 = 0. The equation of l is ( ).}$
\begin{enumerate}
\item $\displaystyle \text{A. 2x - y - 7 = 0}$
\item $\displaystyle \text{B. 2x + y - 5 = 0}$
\item $\displaystyle \text{C. x - 2y - 1 = 0}$
\item $\displaystyle \text{D. x + 2y - 3 = 0}$
\end{enumerate}
\item \textbf{Март, № 45} \par \smallskip
$\displaystyle \text{Suppose that a straight line l is perpendicular to line l}_{1}\text{: x + 2y + 4 = 0, and passes through the intersection of l}_{2}\text{: x + y + 1 = 0 and l}_{3}\text{: 2x + y - 1 = 0. The equation of l is}$
\begin{enumerate}
\item $\displaystyle \text{A. 2x - y - 7 = 0}$
\item $\displaystyle \text{B. x - 2y - 1 = 0}$
\item $\displaystyle \text{C. 2x + y - 5 = 0}$
\item $\displaystyle \text{D. x + 2y - 3 = 0}$
\end{enumerate}
\item \textbf{Январь, № 32} \par \smallskip
$\displaystyle \text{Given two straight lines l}_{1}\text{: ax + (a - 1)y + 3 = 0 and l}_{2}\text{: 2x + ay - 1 = 0. If l}_{1}\text{ ⊥ l}_{2}\text{, then the value of a is ( ).}$
\begin{enumerate}
\item $\displaystyle \text{A. -1}$
\item $\displaystyle \text{B. -1  1}$
\item $\displaystyle \text{C. 1}$
\item $\displaystyle \text{D. 0  -1}$
\end{enumerate}
\end{enumerate}

\subsection*{Тригонометрия — Свойства}
\begin{enumerate}
\item \textbf{Декабрь, № 31} \par \smallskip
$\displaystyle \text{Which of the following statements about the function y = }\cos\text{ x is correct?}$
\begin{enumerate}
\item $\displaystyle \text{A. cos(x + }\pi\text{ ) = }\cos\text{ x}$
\item $\displaystyle \text{B.  / Is an odd function}$
\item $\displaystyle \text{C.  [0, }\pi\text{ ]  / Is monotonically increasing on the interval [0,}\pi\text{ ]}$
\item $\displaystyle \text{D.  2}\pi\text{   / Is periodic with a period of 2}\pi\text{ }$
\end{enumerate}
\item \textbf{Январь, № 30} \par \smallskip
$\displaystyle \text{Given the sine function y = }\sin\text{ x, which of the following statements is incorrect? ( )}$
\begin{enumerate}
\item $\displaystyle \text{A.  1  -1 / ​The max and min values are 1 and -1}$
\item $\displaystyle \text{B.  / It'a s periodic function}$
\item $\displaystyle \text{C.  / It is even}$
\item $\displaystyle \text{D.  }\mathbb{R}\text{ / ​Domain is R​}$
\end{enumerate}
\item \textbf{Январь, № 39} \par \smallskip
$\displaystyle \text{About the tangent function y = }\tan\text{ x, which of the following statements is correct? ( )}$
\begin{enumerate}
\item $\displaystyle \text{A.  / The  range is [-1,1]}$
\item $\displaystyle \text{B.  / The domain is }\mathbb{R}$
\item $\displaystyle \text{C.  / It is even}$
\item $\displaystyle \text{D.  }\pi\text{  / ​Minimum positive period = }\pi\text{ }$
\end{enumerate}
\end{enumerate}

\subsection*{Тригонометрия — Тригонометрические тождества}
\begin{enumerate}
\item \textbf{Март, № 21} \par \smallskip
$\displaystyle \text{Which of the following is correct?}$
\begin{enumerate}
\item $\displaystyle \text{A. }\tan\text{ x = }\cos\text{ x / }\sin\text{ x (}\sin\text{ x }\neq\text{  0)}$
\item $\displaystyle \text{B. }\tan\text{ x = }\sin\text{ x / }\cos\text{ x (}\cos\text{ x }\neq\text{  0)}$
\item $\displaystyle \text{C. }\tan\text{ x = 1 / }\sin\text{ x (}\sin\text{ x }\neq\text{  0)}$
\item $\displaystyle \text{D. }\tan\text{ x = 1 / }\cos\text{ x (}\cos\text{ x }\neq\text{  0)}$
\end{enumerate}
\item \textbf{Март, № 39} \par \smallskip
$\displaystyle \text{If }\sin\text{ }\alpha\text{  = 1/3, then cos(}\pi\text{ /2 + }\alpha\text{ ) =}$
\begin{enumerate}
\item $\displaystyle \text{A. -1/3}$
\item $\displaystyle \text{B. -2}\sqrt{2}\text{/3}$
\item $\displaystyle \text{C. 1/3}$
\item $\displaystyle \text{D. 2}\sqrt{2}\text{/3}$
\end{enumerate}
\item \textbf{Декабрь, № 38} \par \smallskip
$\displaystyle \text{Which of the following derivative formulas is incorrect?}$
\begin{enumerate}
\item $\displaystyle \text{A. cos(2}\pi\text{  - }\alpha\text{ ) = }\cos\text{ }\alpha\text{ }$
\item $\displaystyle \text{B. tan(}\pi\text{  + }\alpha\text{ ) = -}\tan\text{ }\alpha\text{ }$
\item $\displaystyle \text{C. sin(}\pi\text{  - }\alpha\text{ ) = }\sin\text{ }\alpha\text{ }$
\item $\displaystyle \text{D. sin(}\pi\text{ /2 - }\alpha\text{ ) = }\cos\text{ }\alpha\text{ }$
\end{enumerate}
\item \textbf{Январь, № 38} \par \smallskip
$\displaystyle \text{Which of the following statements is correct? ( )}$
\begin{enumerate}
\item $\displaystyle \text{A. sin(}\pi\text{  + }\alpha\text{ ) = }\sin\text{ }\alpha\text{ }$
\item $\displaystyle \text{B. cos(}\pi\text{  + }\alpha\text{ ) = }\cos\text{ }\alpha\text{ }$
\item $\displaystyle \text{C. sin(-}\alpha\text{ ) = }\sin\text{ }\alpha\text{ }$
\item $\displaystyle \text{D. cos(-}\alpha\text{ ) = }\cos\text{ }\alpha\text{ }$
\end{enumerate}
\end{enumerate}

\subsection*{Логарифмические функции — Логарифмы}
\begin{enumerate}
\item \textbf{Март, № 29} \par \smallskip
$\displaystyle \text{Let a = lg 4, b = lg 5. Then a + 2b =}$
\begin{enumerate}
\item $\displaystyle \text{A. 100}$
\item $\displaystyle \text{B. 4}$
\item $\displaystyle \text{C. 5}$
\item $\displaystyle \text{D. 2}$
\end{enumerate}
\item \textbf{Март, № 42} \par \smallskip
$\displaystyle \text{Let f(x) = log}_{5}\text{(2x + 1). Then the solution set of f(x) }\geq\text{  0 is}$
\begin{enumerate}
\item $\displaystyle \text{A. }\{\text{x }\mid\text{ x > 1}\}$
\item $\displaystyle \text{B. }\{\text{x }\mid\text{ x > 0}\}$
\item $\displaystyle \text{C. }\{\text{x }\mid\text{ x }\geq\text{  0}\}$
\item $\displaystyle \text{D. }\{\text{x }\mid\text{ x }\geq\text{  1}\}$
\end{enumerate}
\item \textbf{Январь, № 29} \par \smallskip
$\displaystyle \text{Let a > 0, and b > 0 be real numbers. Then which of the following statements is incorrect? ( )}$
\begin{enumerate}
\item $\displaystyle \text{A. ln(aᵇ) = b }\ln\text{ a}$
\item $\displaystyle \text{B. logₐ b = }\ln\text{ a / }\ln\text{ b (a}\neq\text{ 1, b}\neq\text{ 1)}$
\item $\displaystyle \text{C. ln(a/b) = }\ln\text{ a - }\ln\text{ b}$
\item $\displaystyle \text{D. ln(ab) = }\ln\text{ a + }\ln\text{ b}$
\end{enumerate}
\item \textbf{Декабрь, № 30} \par \smallskip
\textit{(no English statement.)}
\begin{enumerate}
\item $\displaystyle \text{A. -5}$
\item $\displaystyle \text{B. 10}$
\item $\displaystyle \text{C. 0}$
\item $\displaystyle \text{D. 5}$
\end{enumerate}
\item \textbf{Декабрь, № 42} \par \smallskip
$\displaystyle \text{Which of the following statements about the logarithmic function y = logₐx is correct?}$
\begin{enumerate}
\item $\displaystyle \text{A.  (The domain of this function is) }\OPENNEGINF\text{ , }\PLUSINFTAIL\text{ )}$
\item $\displaystyle \text{B.  a > 1  / When a>1, it is a monotonically decreasing function over its domain}$
\item $\displaystyle \text{C.  a > 1  x > 0  logₐ x > 1 / When a>1, for any x>0, we always have logₐx>1}$
\item $\displaystyle \text{D.  0 < a < 1  / When 0<a<1, it is a monotonically decreasing function over its domain}$
\end{enumerate}
\end{enumerate}

\subsection*{Тригонометрия — Синус и косинус суммы}
\begin{enumerate}
\item \textbf{Март, № 35} \par \smallskip
$\displaystyle \text{Let }\alpha\text{ , }\beta\text{  be acute angles. Then which of the following is correct?}$
\begin{enumerate}
\item $\displaystyle \text{A. }\sin\text{ }\alpha\text{  + }\sin\text{ }\beta\text{  < sin(}\alpha\text{  + }\beta\text{ )}$
\item $\displaystyle \text{B. cos(}\alpha\text{  - }\beta\text{ ) < cos(}\alpha\text{  + }\beta\text{ )}$
\item $\displaystyle \text{C. }\cos\text{ }\alpha\text{  }\cos\text{ }\beta\text{  < cos(}\alpha\text{  + }\beta\text{ )}$
\item $\displaystyle \text{D. sin(}\alpha\text{  - }\beta\text{ ) < sin(}\alpha\text{  + }\beta\text{ )}$
\end{enumerate}
\item \textbf{Январь, № 34} \par \smallskip
$\displaystyle \text{Let }\cos\text{ }\alpha\text{  = -}\sqrt{5}\text{/5, and }\alpha\text{  }\in\text{  (}\pi\text{ /2, }\pi\text{ ). Then sin(}\pi\text{ /6 + }\alpha\text{ ) = ( ).}$
\begin{enumerate}
\item $\displaystyle \text{A. (2}\sqrt{5}\text{ - }\sqrt{15}\text{)/10}$
\item $\displaystyle \text{B. (2}\sqrt{15}\text{ + }\sqrt{5}\text{)/10}$
\item $\displaystyle \text{C. (2}\sqrt{15}\text{ - }\sqrt{5}\text{)/10}$
\item $\displaystyle \text{D. (2}\sqrt{5}\text{ + }\sqrt{15}\text{)/10}$
\end{enumerate}
\item \textbf{Декабрь, № 35} \par \smallskip
$\displaystyle \text{Given that }\sin\text{ }\pi\text{ /4 = }\sqrt{2}\text{/2, }\sin\text{ }\pi\text{ /6 = 1/2, then }\cos\text{ }\pi\text{ /12 =}$
\begin{enumerate}
\item $\displaystyle \text{A. }\sqrt{6}\text{/4}$
\item $\displaystyle \text{B. (}\sqrt{6}\text{ + }\sqrt{2}\text{)/4}$
\item $\displaystyle \text{C. (}\sqrt{6}\text{ - }\sqrt{2}\text{)/2}$
\item $\displaystyle \text{D. (1 + }\sqrt{3}\text{)/4}$
\end{enumerate}
\end{enumerate}

\subsection*{Конические сечения — Парабола}
\begin{enumerate}
\item \textbf{Март, № 32} \par \smallskip
$\displaystyle \text{The equation of the directrix of the parabola y = -(1/4)x}^{2}\text{ is}$
\begin{enumerate}
\item $\displaystyle \text{A. y = 1/16}$
\item $\displaystyle \text{B. x = 1/16}$
\item $\displaystyle \text{C. y = 1}$
\item $\displaystyle \text{D. x = 1}$
\end{enumerate}
\item \textbf{Декабрь, № 32} \par \smallskip
$\displaystyle \text{Which of the following statements about the parabola y}^{2}\text{ = 4x is correct?}$
\begin{enumerate}
\item $\displaystyle \text{A.  (Passes through the point)   (-4, 4)}$
\item $\displaystyle \text{B.  (The axis of symmetry is) x = 0}$
\item $\displaystyle \text{C.  (The coordinates of the foci are)  (±1, 0)}$
\item $\displaystyle \text{D.  (The coordinates of the focus are)  (1, 0)}$
\end{enumerate}
\item \textbf{Март, № 37} \par \smallskip
$\displaystyle \text{The equation of the directrix of the parabola y}^{2}\text{ = -x is}$
\begin{enumerate}
\item $\displaystyle \text{A. x = 1/4}$
\item $\displaystyle \text{B. x = 1/2}$
\item $\displaystyle \text{C. x = -1/4}$
\item $\displaystyle \text{D. x = -1/2}$
\end{enumerate}
\item \textbf{Январь, № 31} \par \smallskip
$\displaystyle \text{Let the focus of the parabola y}^{2}\text{ = 4x be F. Suppose that a point P is on the parabola, and the horizontal coordinate of P is 4. Then |PF| = ( ).}$
\begin{enumerate}
\item $\displaystyle \text{A. 2}$
\item $\displaystyle \text{B. 3}$
\item $\displaystyle \text{C. 4}$
\item $\displaystyle \text{D. 5}$
\end{enumerate}
\item \textbf{Декабрь, № 39} \par \smallskip
$\displaystyle \text{The equation of the directrix of the parabola y}^{2}\text{ = -x is}$
\begin{enumerate}
\item $\displaystyle \text{A. x = 1/2}$
\item $\displaystyle \text{B. x = -1/2}$
\item $\displaystyle \text{C. x = 1/4}$
\item $\displaystyle \text{D. x = -1/4}$
\end{enumerate}
\item \textbf{Январь, № 42} \par \smallskip
$\displaystyle \text{The equation of the directrix of the parabola y}^{2}\text{ = -x is ( ).}$
\begin{enumerate}
\item $\displaystyle \text{A. x = 1/4}$
\item $\displaystyle \text{B. x = -1/4}$
\item $\displaystyle \text{C. x = -1/2}$
\item $\displaystyle \text{D. x = 1/2}$
\end{enumerate}
\end{enumerate}

\subsection*{Функции — Тождественные функции}
\begin{enumerate}
\item \textbf{Декабрь, № 36} \par \smallskip
$\displaystyle \text{Which of the following functions is identical to y = |x|?}$
\begin{enumerate}
\item $\displaystyle \text{A. y = ³√x}^{3}$
\item $\displaystyle \text{B. y = (√x)}^{2}$
\item $\displaystyle \text{C. y = x}$
\item $\displaystyle \text{D. y = √x}^{2}$
\end{enumerate}
\item \textbf{Январь, № 35} \par \smallskip
$\displaystyle \text{Which of the following are the same function ? ( )}$
\begin{enumerate}
\item $\displaystyle \text{A. y = (√x)}^{2}\text{  and y = }\sqrt{x^{2}}$
\item $\displaystyle \text{B. y = (x}^{4}\text{ - 1)/(x}^{2}\text{ + 1)   and y = x}^{2}\text{ - 1}$
\item $\displaystyle \text{C. y = 1  and y = x}^{2}$
\item $\displaystyle \text{D. y = x  and y = }\sqrt{x^{2}}$
\end{enumerate}
\end{enumerate}

\subsection*{Тригонометрия — Формулы половинного угла}
\begin{enumerate}
\item \textbf{Март, № 34} \par \smallskip
$\displaystyle \text{Suppose that }\cos\text{ }\alpha\text{  = -1/2, }\alpha\text{  }\in\text{  (}\pi\text{ /2, }\pi\text{ ), then the value of sin(}\alpha\text{ /2) is}$
\begin{enumerate}
\item $\displaystyle \text{A. }\sqrt{3}\text{/2}$
\item $\displaystyle \text{B. -}\sqrt{3}\text{/2}$
\item $\displaystyle \text{C. -1/2}$
\item $\displaystyle \text{D. 1/2}$
\end{enumerate}
\item \textbf{Декабрь, № 37} \par \smallskip
$\displaystyle \text{If }\alpha\text{  is an acute angle, and }\cos\text{ }\alpha\text{  = 1/2, then }\sin\text{ }\alpha\text{ /2 =}$
\begin{enumerate}
\item $\displaystyle \text{A. ±}\sqrt{3}\text{/2}$
\item $\displaystyle \text{B. ±1/2}$
\item $\displaystyle \text{C. }\sqrt{3}\text{/2}$
\item $\displaystyle \text{D. 1/2}$
\end{enumerate}
\item \textbf{Январь, № 37} \par \smallskip
$\displaystyle \text{Suppose that }\cos\text{ }\alpha\text{  = -1/2, }\alpha\text{  }\in\text{  (}\pi\text{ /2, }\pi\text{ ), then the value of sin(}\alpha\text{ /2) is ( ).}$
\begin{enumerate}
\item $\displaystyle \text{A. -}\sqrt{3}\text{/2}$
\item $\displaystyle \text{B. }\sqrt{3}\text{/2}$
\item $\displaystyle \text{C. -1/2}$
\item $\displaystyle \text{D. 1/2}$
\end{enumerate}
\item \textbf{Январь, № 36} \par \smallskip
$\displaystyle \text{Let }\sin\text{ }\alpha\text{  = -12/13, and }\alpha\text{  }\in\text{  (}\pi\text{ , 3}\pi\text{ /2). Then cos(}\alpha\text{ /2) = ( ).}$
\begin{enumerate}
\item $\displaystyle \text{A. 3}\sqrt{13}\text{/13}$
\item $\displaystyle \text{B. -3}\sqrt{13}\text{/13}$
\item $\displaystyle \text{C. -2}\sqrt{13}\text{/13}$
\item $\displaystyle \text{D. 2}\sqrt{13}\text{/13}$
\end{enumerate}
\item \textbf{Март, № 38} \par \smallskip
$\displaystyle \text{If }\cos\text{ }\alpha\text{  = 1/3 and -}\pi\text{ /2 < }\alpha\text{  < 0, then tan(}\alpha\text{ /2) =}$
\begin{enumerate}
\item $\displaystyle \text{A. -}\sqrt{2}$
\item $\displaystyle \text{B. }\sqrt{2}\text{/2}$
\item $\displaystyle \text{C. -}\sqrt{2}\text{/2}$
\item $\displaystyle \text{D. }\sqrt{2}$
\end{enumerate}
\end{enumerate}

\subsection*{Функции — Графики}
\begin{enumerate}
\item \textbf{Январь, № 41} \par \smallskip
$\displaystyle \text{If the graph of the function f(x) = 2 + logₐ(x - 3) (a > 0, and a }\neq\text{  1) passes through a fixed point P, then the coordinates of P are ( ).}$
\begin{enumerate}
\item $\displaystyle \text{A. (3,1)}$
\item $\displaystyle \text{B. (3,0)}$
\item $\displaystyle \text{C. (4,0)}$
\item $\displaystyle \text{D. (4,2)}$
\end{enumerate}
\end{enumerate}

\subsection*{Конические сечения — Эллипс}
\begin{enumerate}
\item \textbf{Март, № 43} \par \smallskip
$\displaystyle \text{Suppose that the major axis of an ellipse is 10, the two foci F}_{1}\text{, F}_{2}\text{ are located at (3,0) and (-3,0). The equation of the ellipse is}$
\begin{enumerate}
\item $\displaystyle \text{A. x}^{2}\text{/25 + y}^{2}\text{/16 = 1}$
\item $\displaystyle \text{B. x}^{2}\text{/16 + y}^{2}\text{/9 = 1}$
\item $\displaystyle \text{C. x}^{2}\text{/9 + y}^{2}\text{/25 = 1}$
\item $\displaystyle \text{D. x}^{2}\text{/16 + y}^{2}\text{/25 = 1}$
\end{enumerate}
\item \textbf{Январь, № 43} \par \smallskip
$\displaystyle \text{Suppose that the foci of the ellipse x}^{2}\text{/4 + y}^{2}\text{/2 = 1 are F}_{1}\text{, F}_{2}\text{, and the point P is on the ellipse. Then |PF}_{1}\text{| + |PF}_{2}\text{| = ( ).}$
\begin{enumerate}
\item $\displaystyle \text{A. 2}$
\item $\displaystyle \text{B. 4}$
\item $\displaystyle \text{C. 8}$
\item $\displaystyle \text{D. 6}$
\end{enumerate}
\item \textbf{Декабрь, № 43} \par \smallskip
$\displaystyle \text{Which of the following statements about the ellipse x}^{2}\text{/4 + y}^{2}\text{/20 = 1 is correct?}$
\begin{enumerate}
\item $\displaystyle \text{A.  (The foci of this ellipse are) (0, ±4)}$
\item $\displaystyle \text{B.  (The foci of this ellipse are) (±4, 0)}$
\item $\displaystyle \text{C.  (The vertices of this ellipse are ) (±4, 0)  (and) (0, ±20)}$
\item $\displaystyle \text{D.   (The lengths of the major and minor axes of this ellipse are ) 20  (and) 4}$
\end{enumerate}
\end{enumerate}

\subsection*{Геометрия — Векторы}
\begin{enumerate}
\item \textbf{Март, № 44} \par \smallskip
$\displaystyle \text{Let a = (-2,3), b = (2,1). Then a - 2b =}$
\begin{enumerate}
\item $\displaystyle \text{A. (-6,1)}$
\item $\displaystyle \text{B. (-6,-1)}$
\item $\displaystyle \text{C. (6,-1)}$
\item $\displaystyle \text{D. (6,1)}$
\end{enumerate}
\item \textbf{Декабрь, № 45} \par \smallskip
$\displaystyle \text{Given that O is the origin, A and B are two points on the plane, and P is the midpoint of AB, then the vector OP=}$
\begin{enumerate}
\item $\displaystyle \text{A. 1/2(OA + OB)}$
\item $\displaystyle \text{B. OA + OB}$
\item $\displaystyle \text{C. OB - OA}$
\item $\displaystyle \text{D. 1/2(OA - OB)}$
\end{enumerate}
\item \textbf{Январь, № 44} \par \smallskip
$\displaystyle \text{Given that vector AB = a + 5b, BC = -2a + 8b, CD = 3a - 3b. Thus, ( ).}$
\begin{enumerate}
\item $\displaystyle \text{A. A, C, D  аre on the same line}$
\item $\displaystyle \text{B. A, B, D  аre on the same line}$
\item $\displaystyle \text{C. B, C, D  аre on the same line}$
\item $\displaystyle \text{D. A, B, C  аre on the same line}$
\end{enumerate}
\end{enumerate}

\subsection*{Последовательности — Сложные задачи}
\begin{enumerate}
\item \textbf{Декабрь, № 47} \par \smallskip
$\displaystyle \text{Given that the arithmetic sequence \{a}_{n}\text{\} satisfies 3(a}_{2}\text{ + a}_{6}\text{) + 2(a}_{6}\text{ + a}_{10}\text{ + a}_{14}\text{) = 24, then S}_{13}\text{ =}$
\begin{enumerate}
\item $\displaystyle \text{A. 104}$
\item $\displaystyle \text{B. 52}$
\item $\displaystyle \text{C. 39}$
\item $\displaystyle \text{D. 26}$
\end{enumerate}
\item \textbf{Декабрь, № 41} \par \smallskip
$\displaystyle \text{Given that the sequence \{a}_{n}\text{\} satisfies a}_{1}\text{ = 1, a}_{n}\text{ = 1/(1 + 1/a}_{n}\text{₋₁) (n }\geq\text{  2), then a}_{100}\text{ =}$
\begin{enumerate}
\item $\displaystyle \text{A. 50}$
\item $\displaystyle \text{B. 1/100}$
\item $\displaystyle \text{C. 1}$
\item $\displaystyle \text{D. 100}$
\end{enumerate}
\item \textbf{Январь, № 47} \par \smallskip
$\displaystyle \text{Suppose that, in a sequence \{a}_{n}\text{\}, we have a}_{1}\text{ = 1, and 1/a}_{n}\text{ + 2/a}_{n}\text{₊₁ = 0 (n = 1,2,}\ldots\text{ ). Suppose that another sequence \{b}_{n}\text{\} satisfies b}_{n}\text{ = |a}_{n}\text{| (n = 1,2,}\ldots\text{ ). Then, the sum of the first n terms of the sequence \{b}_{n}\text{\} is S}_{n}\text{ = ( ).}$
\begin{enumerate}
\item $\displaystyle \text{A. (1 + 2}^{n}\text{)/3}$
\item $\displaystyle \text{B. |1 - (-2)}^{n}\text{|/3}$
\item $\displaystyle \text{C. 2}^{n}\text{ - 1}$
\item $\displaystyle \text{D. (-2)}^{n}\text{ - 1}$
\end{enumerate}
\end{enumerate}

\subsection*{Комплексные числа — Сложные задачи}
\begin{enumerate}
\item \textbf{Декабрь, № 46} \par \smallskip
$\displaystyle \text{If the complex number z satisfies z}^{3}\text{ = 1 and z }\neq\text{  1, then 1 + z + z}^{2}\text{ + ... + z}^{999}\text{ =}$
\begin{enumerate}
\item $\displaystyle \text{A. (1 ± }\sqrt{3}\text{i)/2}$
\item $\displaystyle \text{B. (-1 ±}\sqrt{3}\text{i)/2}$
\item $\displaystyle \text{C. 1}$
\item $\displaystyle \text{D. i}$
\end{enumerate}
\item \textbf{Январь, № 46} \par \smallskip
$\displaystyle \text{On the complex plane, the point representing a complex number z is on the straight line x - y = 0. If z is a root of a quadratic equation x}^{2}\text{ + mx + 4 = 0, then m = ( ).}$
\begin{enumerate}
\item $\displaystyle \text{A. }\sqrt{2}\text{  or 2}\sqrt{2}$
\item $\displaystyle \text{B. 2}\sqrt{2}\text{  or -2}\sqrt{2}$
\item $\displaystyle \text{C. -}\sqrt{2}\text{  or -2}\sqrt{2}$
\item $\displaystyle \text{D. }\sqrt{2}\text{  or -}\sqrt{2}$
\end{enumerate}
\item \textbf{Март, № 46} \par \smallskip
$\displaystyle \text{Suppose that z}_{1}\text{ and z}_{2}\text{ are two roots of the equation z}^{2}\text{ - 2z + 2 = 0. Let a}_{n}\text{ = (z}_{1}\text{ⁿ - z}_{2}\text{ⁿ)}^{2}\text{. Then a}_{4}\text{ =}$
\begin{enumerate}
\item $\displaystyle \text{A. 32}$
\item $\displaystyle \text{B. -64}$
\item $\displaystyle \text{C. 0}$
\item $\displaystyle \text{D. -32}$
\end{enumerate}
\end{enumerate}

\subsection*{Вероятность — Элементарная вероятность}
\begin{enumerate}
\item \textbf{Март, № 48} \par \smallskip
$\displaystyle \text{A bag contains five balls identical in size and shape, including two red balls, labeled 1 and 2, two black balls, labeled 1 and 2, one white ball, labeled 1. Two balls are drawn randomly without replacement. The probability that the balls have different colors, and the probability that at least one of the two balls is labeled 1 are respectively}$
\begin{enumerate}
\item $\displaystyle \text{A. 4/5, 7/10}$
\item $\displaystyle \text{B. 4/5, 9/10}$
\item $\displaystyle \text{C. 3/5, 7/10}$
\item $\displaystyle \text{D. 3/5, 9/10}$
\end{enumerate}
\item \textbf{Январь, № 48} \par \smallskip
$\displaystyle \text{In daily life, people enjoy using running apps to track and share their exercise routes. To understand the usage of running apps among high school and college students in a certain region, 200 high school students and 80 college students were randomly chosen, and their favorite running apps were recorded as follows. Assuming that college students' and high school students' preferences for running apps are independent of each other, if one person is randomly selected from each group, the probability that their favorite running apps are different is ( ).}$
\begin{enumerate}
\item $\displaystyle \text{A. 23/80}$
\item $\displaystyle \text{B. 57/80}$
\item $\displaystyle \text{C. 9/20}$
\item $\displaystyle \text{D. 11/20}$
\end{enumerate}
\item \textbf{Декабрь, № 48} \par \smallskip
$\displaystyle \text{The twelve zodiac signs, also known as 'Sheng Xiao', are twelve animals used in China to represent years, in order: Rat, Ox, Tiger, Rabbit, Dragon, Snake, Horse, Sheep, Monkey, Rooster, Dog, Pig. For example, people born in 1984, 1996, and 2008 are of the Rat sign, while those born in 1985, 1997, and 2009 are of the Ox sign, and so on. The probability that at least two out of three randomly selected people have the same zodiac sign is}$
\begin{enumerate}
\item $\displaystyle \text{A. 17/72}$
\item $\displaystyle \text{B. 1/12}$
\item $\displaystyle \text{C. 55/72}$
\item $\displaystyle \text{D. 143/144}$
\end{enumerate}
\end{enumerate}

\clearpage
\section*{Часть 2. Ответы}

\subsection*{Множества — Операции над множествами}
\begin{itemize}
\item Декабрь, № 1 — C
\item Январь, № 1 — A
\item Март, № 1 — B
\item Декабрь, № 2 — B
\item Январь, № 2 — C
\item Март, № 2 — A
\end{itemize}

\subsection*{Неравенства — Квадратичные неравенства}
\begin{itemize}
\item Март, № 3 — A
\item Декабрь, № 3 — D
\item Январь, № 3 — A
\end{itemize}

\subsection*{Функции — Область определения}
\begin{itemize}
\item Март, № 4 — D
\item Декабрь, № 4 — B
\item Январь, № 5 — C
\item Декабрь, № 20 — A
\end{itemize}

\subsection*{Функции — Обратные функции}
\begin{itemize}
\item Январь, № 9 — A
\item Декабрь, № 9 — D
\item Март, № 8 — D
\end{itemize}

\subsection*{Геометрия — Координатная геометрия}
\begin{itemize}
\item Январь, № 10 — B
\item Декабрь, № 11 — A
\item Март, № 5 — C
\item Декабрь, № 6 — A
\item Январь, № 6 — B
\item Март, № 16 — C
\end{itemize}

\subsection*{Неравенства — Рациональные неравенства}
\begin{itemize}
\item Декабрь, № 12 — C
\item Март, № 11 — C
\item Январь, № 12 — A
\end{itemize}

\subsection*{Среднее арифметическое и геометрическое — Среднее геометрическое}
\begin{itemize}
\item Март, № 12 — D
\end{itemize}

\subsection*{Среднее арифметическое и геометрическое — Среднее арифметическое}
\begin{itemize}
\item Январь, № 13 — B
\end{itemize}

\subsection*{Функции — Равенство функций}
\begin{itemize}
\item Декабрь, № 15 — D
\end{itemize}

\subsection*{Тригонометрия — Сторона угла через точку}
\begin{itemize}
\item Январь, № 15 — B
\item Декабрь, № 16 — D
\end{itemize}

\subsection*{Тригонометрия — Табличные значения}
\begin{itemize}
\item Январь, № 7 — B
\item Декабрь, № 7 — C
\item Декабрь, № 22 — B
\item Март, № 6 — A
\item Январь, № 20 — B
\item Декабрь, № 34 — B
\item Декабрь, № 19 — B
\end{itemize}

\subsection*{Функции — Свойства функций}
\begin{itemize}
\item Март, № 7 — D
\item Январь, № 8 — D
\item Март, № 9 — A
\item Январь, № 11 — B
\item Март, № 31 — C
\item Декабрь, № 8 — D
\item Декабрь, № 10 — C
\item Январь, № 23 — C
\item Декабрь, № 26 — D
\item Март, № 36 — A
\item Март, № 25 — A
\item Март, № 40 — B
\end{itemize}

\subsection*{Последовательности — Простые задачи}
\begin{itemize}
\item Март, № 10 — B
\item Март, № 22 — A
\item Январь, № 4 — A
\item Декабрь, № 5 — A
\item Декабрь, № 13 — A
\item Январь, № 16 — A
\item Декабрь, № 17 — B
\item Январь, № 21 — A
\item Март, № 15 — B
\item Декабрь, № 23 — C
\item Январь, № 40 — B
\item Март, № 41 — A
\item Март, № 47 — D
\end{itemize}

\subsection*{Неравенства — Свойства неравенств}
\begin{itemize}
\item Январь, № 22 — B
\item Март, № 23 — A
\item Март, № 19 — A
\end{itemize}

\subsection*{Конические сечения — Окружность}
\begin{itemize}
\item Март, № 20 — B
\item Март, № 24 — B
\item Январь, № 28 — D
\item Декабрь, № 18 — D
\item Январь, № 19 — C
\item Декабрь, № 21 — C
\end{itemize}

\subsection*{Геометрия — Формула расстояния}
\begin{itemize}
\item Январь, № 17 — B
\item Декабрь, № 25 — A
\item Март, № 30 — D
\item Январь, № 33 — B
\item Март, № 17 — D
\end{itemize}

\subsection*{Функции — Неравенства}
\begin{itemize}
\item Декабрь, № 24 — C
\item Январь, № 27 — D
\end{itemize}

\subsection*{Тригонометрия — Формулы двойного угла}
\begin{itemize}
\item Март, № 26 — B
\item Декабрь, № 27 — C
\item Январь, № 24 — B
\end{itemize}

\subsection*{Тригонометрия — Тригонометрические выражения}
\begin{itemize}
\item Март, № 14 — A
\item Декабрь, № 40 — A
\end{itemize}

\subsection*{Конические сечения — Гипербола}
\begin{itemize}
\item Март, № 28 — B
\item Декабрь, № 29 — A
\item Январь, № 26 — A
\end{itemize}

\subsection*{Геометрия — Прямые}
\begin{itemize}
\item Март, № 13 — A
\item Декабрь, № 14 — A
\item Март, № 18 — D
\item Март, № 33 — B
\item Январь, № 14 — C
\item Январь, № 18 — B
\item Январь, № 25 — D
\item Март, № 27 — B
\item Декабрь, № 28 — D
\item Декабрь, № 33 — B
\item Декабрь, № 44 — C
\item Январь, № 45 — A
\item Март, № 45 — A
\item Январь, № 32 — D
\end{itemize}

\subsection*{Тригонометрия — Свойства}
\begin{itemize}
\item Декабрь, № 31 — D
\item Январь, № 30 — C
\item Январь, № 39 — D
\end{itemize}

\subsection*{Тригонометрия — Тригонометрические тождества}
\begin{itemize}
\item Март, № 21 — B
\item Март, № 39 — A
\item Декабрь, № 38 — B
\item Январь, № 38 — D
\end{itemize}

\subsection*{Логарифмические функции — Логарифмы}
\begin{itemize}
\item Март, № 29 — D
\item Март, № 42 — C
\item Январь, № 29 — B
\item Декабрь, № 30 — C
\item Декабрь, № 42 — D
\end{itemize}

\subsection*{Тригонометрия — Синус и косинус суммы}
\begin{itemize}
\item Март, № 35 — D
\item Январь, № 34 — C
\item Декабрь, № 35 — B
\end{itemize}

\subsection*{Конические сечения — Парабола}
\begin{itemize}
\item Март, № 32 — С
\item Декабрь, № 32 — D
\item Март, № 37 — A
\item Январь, № 31 — D
\item Декабрь, № 39 — C
\item Январь, № 42 — A
\end{itemize}

\subsection*{Функции — Тождественные функции}
\begin{itemize}
\item Декабрь, № 36 — D
\item Январь, № 35 — B
\end{itemize}

\subsection*{Тригонометрия — Формулы половинного угла}
\begin{itemize}
\item Март, № 34 — A
\item Декабрь, № 37 — D
\item Январь, № 37 — B
\item Январь, № 36 — C
\item Март, № 38 — C
\end{itemize}

\subsection*{Функции — Графики}
\begin{itemize}
\item Январь, № 41 — D
\end{itemize}

\subsection*{Конические сечения — Эллипс}
\begin{itemize}
\item Март, № 43 — A
\item Январь, № 43 — B
\item Декабрь, № 43 — A
\end{itemize}

\subsection*{Геометрия — Векторы}
\begin{itemize}
\item Март, № 44 — A
\item Декабрь, № 45 — A
\item Январь, № 44 — B
\end{itemize}

\subsection*{Последовательности — Сложные задачи}
\begin{itemize}
\item Декабрь, № 47 — D
\item Декабрь, № 41 — B
\item Январь, № 47 — C
\end{itemize}

\subsection*{Комплексные числа — Сложные задачи}
\begin{itemize}
\item Декабрь, № 46 — C
\item Январь, № 46 — B
\item Март, № 46 — C
\end{itemize}

\subsection*{Вероятность — Элементарная вероятность}
\begin{itemize}
\item Март, № 48 — B
\item Январь, № 48 — B
\item Декабрь, № 48 — A
\end{itemize}

\end{document}